% Cote 119: offprints, reports, courses, CV and letters (typed and printed part): archive cote 119, site page 187 (PDF page 188).
% typescript. Draft compared: none.
% AI-transcribed by Claude Opus 5.5 (claude-opus-5-5, adaptive thinking, effort high), Message Batches API, batch r119_full (msgbatch_01DwZNm5Y7YsKBWyAN4sFs8w), 2026-11-13; not yet checked by eye.
\scanpage{119/187}
\mspage{- 6 -}

La fonction $g(u) = \Sigma_{n > 0}\ a_{-n}\ u^{n}$ est holomorphe pour $|u| < (\rho_2)^{-1}$.

La fonction $z \longrightarrow u = \frac{1}{z}$ est aussi holomorphe dans cet ouvert, donc $f_2$ est
holomorphe comme composée de deux fonctions holomorphes.

Si $\rho_1 > \rho_2$ la somme $f(z)$ de la série $\Sigma_{n \in \mathbb{Z}}\ a_n\ z^{n}$ est donc ho-%
lomorphe dans la couronne circulaire $f_2 < |z| < \rho_1$ et la convergence de la
série est normale dans tout couronne $r_2 \leq |z| \leq r_1$ pour tous $r_1$ et $r_2$ tels
que $\rho_2 < r_2 < r_1 < \rho_1$ .
La dérivée de $f_2(z) = g(\frac{1}{z})$ est égale à $f_2'(z) = - \frac{1}{z^2}\, g'(\frac{1}{z}) = \Sigma_{n < 0}\ n\, a_n\ z^{n-1}$

car on a : \quad $g'(u) = \Sigma_{n > 0}\ n\, a_n\ u^{n-1}$ .

La dérivée de $f(z)$ est donc obtenue en dérivant terme à terme la série de
Laurent $\Sigma_{n \in z}\ a_n\ z^{n}$.

N.B. On admet dans ce qui précède les cas $\rho_2 = 0$ $(\frac{1}{\rho_2} = + \infty)$ et $\rho_1 = + \infty$.
8.2.2. \uline{Théorème}.

Toute fonction holomorphe $f$ dans une couronne $\rho_2 < |z| < \rho_1$ est
développable en série de Laurent dans cette couronne.
Autrement dit il existe une série de Laurent $\Sigma\, a_n\ z^{n}$ qui converge dans cette
couronne et dont la somme est égale à $f(z)$ en tout point de la couronne.
\uline{Démonstration} :

Soit $r_1$ et $r_2$ tels que $\rho_2 < r_2 < r_1 < \rho_1$ . Si la fonction $f$ est
développable en série de Laurent dans la couronne $\rho_2 < |z| < \rho_1$ et si l'on pose
$z = re^{i\theta}$ d'après 8.2.1 la convergence de la série est normale dans toute couron-%
ne $r_2 \leq |z| \leq r_1$ et donc sur tout cercle $|z| = r$ pour $\rho_2 < r < \rho_1$ .

A partir de $f(re^{i\theta}) = \Sigma_{n \in \mathbb{Z}}\ a_n\ r^{n}\ e^{in\theta}$

on obtient donc comme dans (7.1)
\[
(8.2.3) \qquad a_n\ r^{n} = \frac{1}{2\pi} \int_0^{2\pi} e^{-in\theta}\ f(r\, e^{i\theta})\ d\theta \qquad (n \in \mathbb{Z})
\]

\margin{187}
